% the page size

\documentclass{article}%
\usepackage{hyperref}
\usepackage{color}
\usepackage{amsmath}
\usepackage{amsfonts}
\usepackage{amssymb}
\usepackage{graphicx}
\usepackage{geometry}%
\setcounter{MaxMatrixCols}{30}
%TCIDATA{OutputFilter=latex2.dll}
%TCIDATA{Version=5.50.0.2953}
%TCIDATA{CSTFile=40 LaTeX article.cst}
%TCIDATA{Created=Thursday, October 06, 2016 16:43:40}
%TCIDATA{LastRevised=Friday, October 12, 2018 08:59:35}
%TCIDATA{<META NAME="GraphicsSave" CONTENT="32">}
%TCIDATA{<META NAME="SaveForMode" CONTENT="1">}
%TCIDATA{BibliographyScheme=Manual}
%TCIDATA{<META NAME="DocumentShell" CONTENT="Standard LaTeX\Blank - Standard LaTeX Article">}
%TCIDATA{Language=American English}
%BeginMSIPreambleData
\providecommand{\U}[1]{\protect\rule{.1in}{.1in}}
%EndMSIPreambleData
\geometry{a4paper, top=1cm, left=1cm, right=1cm, bottom=1.0cm, includehead, includefoot}
\graphicspath{{C:/Users/Gauthier/Dropbox/TD/2015_DSGE/chap3/codes/}}
\newtheorem{theorem}{Theorem}
\newtheorem{acknowledgement}[theorem]{Acknowledgement}
\newtheorem{algorithm}[theorem]{Algorithm}
\newtheorem{axiom}[theorem]{Axiom}
\newtheorem{case}[theorem]{Case}
\newtheorem{claim}[theorem]{Claim}
\newtheorem{conclusion}[theorem]{Conclusion}
\newtheorem{condition}[theorem]{Condition}
\newtheorem{conjecture}[theorem]{Conjecture}
\newtheorem{corollary}[theorem]{Corollary}
\newtheorem{criterion}[theorem]{Criterion}
\newtheorem{definition}[theorem]{Definition}
\newtheorem{example}[theorem]{Example}
\newtheorem{exercise}[theorem]{Exercise}
\newtheorem{lemma}[theorem]{Lemma}
\newtheorem{notation}[theorem]{Notation}
\newtheorem{problem}[theorem]{Problem}
\newtheorem{proposition}[theorem]{Proposition}
\newtheorem{remark}[theorem]{Remark}
\newtheorem{solution}[theorem]{Solution}
\newtheorem{summary}[theorem]{Summary}
\newenvironment{proof}[1][Proof]{\noindent\textbf{#1.} }{\ \rule{0.5em}{0.5em}}
\begin{document}

\title{A note on variable capital utilization}
\author{G.\ Vermandel}
\maketitle

\section{Model}

Variable capital utilization is an important pattern of business cycles
fluctuations as underlined by \autoref{fig}

\begin{figure}[ph]
\begin{center}
\includegraphics[
width=\textwidth
]{exo1_u_rate.pdf}
\end{center}
\par
\label{fig}\caption{Capital utilization rate - US economy}%
\end{figure}

Following Christiano et al.\ (2005), it is also possible to introduce variable
capital utilisation in the production technology of firms:%
\begin{equation}
Y_{it}=\varepsilon_{t}^{A}\left(  K_{it}^{u}{}\right)  ^{\alpha}H_{it}%
{}^{1-\alpha},
\end{equation}
such that $K_{it}^{u}$\ denotes the capital utilised by the firms:
\begin{equation}
K_{it}^{u}=u_{t}K_{it-1},
\end{equation}
where $u_{t}$\ denotes the utilised level of capital with $\bar{u}=1$,
$K_{it}$\ is the amount of physical employed in the production process of
goods. The utilisation rate may vary over the cycles: if $u_{t}>1$, capital is
over-utilised, which will force households to increase the capital supply,
while if $u_{t}<1$, capital is under-utilised.\newline Turning to households,
they choose the optimal level of capital utilisation, $u_{jt}$.
Over-utilisation of capital is costly for households in a proportion
$\Phi\left(  u_{jt}\right)  $:%
\begin{equation}
\Phi\left(  u_{jt}\right)  =\bar{Z}\left(  u_{jt}-1\right)  +\frac{1}{2}%
\frac{\psi}{1-\psi}\bar{Z}\left(  u_{jt}-1\right)  ^{2}%
\end{equation}
where $\bar{Z}$\ is the steady state cost of capital services and $\psi$ is
the elasticity of capital utilisation. The budget constraint of households is
now:
\begin{equation}
W_{t}H_{jt}+Z_{t}u_{jt}K_{jt-1}+R_{t-1}B_{jt-1}=C_{jt}+I_{jt}+\Phi\left(
u_{t}\right)  K_{t-1}+B_{jt}+T_{jt}, \label{eq budget}%
\end{equation}
Finally, the resource constraint is also affected by the utilisation cost paid
in terms of goods:%
\begin{equation}
Y_{t}=C_{t}+I_{t}+\bar{G}\varepsilon_{t}^{G}+\Phi\left(  u_{jt}\right)
K_{jt-1}%
\end{equation}


\subsection{Households}

Since \autoref{eq budget} is affected by capital utilization.\ We have to
recalculate the first order conditions of the household optimization scheme.
In addition, the capital utilization rate is optimally determined by
households as they are the suppliers of physical capital to firms.\ 

The Lagrangian problem is given by:%

\begin{align*}
\mathcal{L}_{t}  &  =\mathbb{E}_{t}%
%TCIMACRO{\dsum \limits_{\tau=0}^{\infty}}%
%BeginExpansion
{\displaystyle\sum\limits_{\tau=0}^{\infty}}
%EndExpansion
\varepsilon_{t+\tau}^{C}\beta^{\tau}\{\frac{C_{jt+\tau}^{1-\sigma_{C}}%
}{1-\sigma_{C}}-\chi\frac{H_{jt+\tau}{}^{1+\sigma_{L}}}{1+\sigma_{L}}\\
&  +\lambda_{t+\tau}^{c}\left[  W_{t+\tau}H_{jt+\tau}+Z_{t+\tau}%
\textcolor{red}{u_{jt+\tau}}K_{jt-1+\tau}+R_{t-1+\tau}B_{jt-1+\tau}%
-C_{jt+\tau}-I_{jt+\tau}-B_{jt+\tau}%
-\textcolor{red}{\Phi(u_{jt+\tau})K_{jt-1+\tau}}-T_{jt+\tau}\right] \\
&  +\lambda_{t+\tau}^{k}\left[  \varepsilon_{t+\tau}^{I}I_{jt+\tau}+\left(
1-\delta\right)  K_{jt-1+\tau}-K_{jt+\tau}\right]  \}
\end{align*}


If you're not comfortable with infinite sums, you can develop the sum for
$\tau=0$ and $\tau=1$ (we don't mind for value of tau above 1):%

\begin{align*}
\mathcal{L}_{t}  &  =\varepsilon_{t}^{C}\left(  \frac{C_{jt}^{1-\sigma_{C}}%
}{1-\sigma_{C}}-\chi\frac{H_{jt}{}^{1+\sigma_{L}}}{1+\sigma_{L}}\right) \\
&  +\varepsilon_{t}^{C}\lambda_{t}^{c}\left[  W_{t}H_{jt}+Z_{t}u_{jt}%
K_{jt-1}+R_{t-1}B_{jt-1}-C_{jt}-I_{jt}-B_{jt}-\Phi\left(  u_{jt}\right)
K_{jt-1}-T_{jt}\right] \\
&  +\varepsilon_{t}^{C}\lambda_{t}^{k}\left[  \varepsilon_{t}^{I}%
I_{jt}+\left(  1-\delta\right)  K_{jt-1}-K_{jt}\right]  \}\\
&  +\varepsilon_{t+1}^{C}\beta\left(  \frac{C_{jt+1}^{1-\sigma_{C}}}%
{1-\sigma_{C}}-\chi\frac{H_{jt+1}{}^{1+\sigma_{L}}}{1+\sigma_{L}}\right) \\
&  +\varepsilon_{t+1}^{C}\lambda_{t+1}^{c}\beta\left[  W_{t+1}H_{jt+1}%
+Z_{t+1}u_{jt+1}K_{jt}+R_{t}B_{jt}-C_{jt+1}-I_{jt+1}-B_{jt+1}-\Phi\left(
u_{jt+1}\right)  K_{jt}-T_{jt+1}\right] \\
&  +\varepsilon_{t+1}^{C}\lambda_{t+1}^{k}\beta\left[  \varepsilon_{t}%
^{I}I_{jt+1}+\left(  1-\delta\right)  K_{jt}-K_{jt+1}\right] \\
&  +\varepsilon_{t+2}^{C}\beta^{2}\left(  \frac{C_{jt+2}^{1-\sigma_{C}}%
}{1-\sigma_{C}}-\chi\frac{H_{jt+2}{}^{1+\sigma_{L}}}{1+\sigma_{L}}\right) \\
&  +....
\end{align*}


Only two first order conditions are different from what we had in the standard
RBC model:%
\begin{align*}
\left(  \partial C_{jt}\right)   &  :\lambda_{t}^{c}=C_{jt}{}^{-\sigma_{C}}\\
\left(  \partial H_{jt}\right)   &  :\chi H_{jt}^{\sigma^{L}}=\lambda_{t}%
^{c}W_{t}\\
\left(  \partial B_{jt}\right)   &  :\mathbb{E}_{t}\left\{  \frac
{\varepsilon_{t}^{C}\lambda_{t}^{c}}{\varepsilon_{t+1}^{C}\lambda_{t+1}^{c}%
}\right\}  =\beta R_{t}\\
\left(  \partial I_{jt}\right)   &  :\lambda_{t}^{c}=\varepsilon_{t}%
^{I}\lambda_{t}^{k}\\
\left(  \partial K_{jt}\right)   &  :\varepsilon_{t}^{C}\lambda_{t}^{k}%
=\beta\mathbb{E}_{t}\left\{  \varepsilon_{t+1}^{C}\left[  \overset{}{\left(
1-\delta\right)  }\lambda_{t+1}^{k}+\lambda_{t+1}^{c}%
\textcolor{red}{u_{jt+1}}Z_{t+1}-\lambda_{t+1}^{c}%
\textcolor{red}{\Phi(u_{jt+1})}\right]  \right\}  \\
&
\textcolor{red}{(\partial u_{jt}):\epsilon_{t}^{C}\lambda_{t}^{c}[Z_{t}K_{jt-1}-\Phi'(u_{jt})K_{jt-1}]=0}
\end{align*}


We can substitute the Lagrangian multipliers in $\left(  \partial
K_{jt}\right)  $ by injecting $\left(  \partial I_{jt}\right)  $ and replace
$\lambda_{t}^{k}=\lambda_{t}^{c}/\varepsilon_{t}^{I}$:\
\begin{align*}
\left(  \partial K_{jt}\right)   &  :\varepsilon_{t}^{C}\lambda_{t}^{k}%
=\beta\mathbb{E}_{t}\left\{  \varepsilon_{t+1}^{C}\left[  \overset{}{\left(
1-\delta\right)  }\lambda_{t+1}^{k}+\lambda_{t+1}^{c}%
\textcolor{red}{u_{jt+1}}Z_{t+1}-\lambda_{t+1}^{c}%
\textcolor{red}{\Phi(u_{jt+1})}\right]  \right\} \\
&  :\varepsilon_{t}^{C}\frac{\lambda_{t}^{c}}{\varepsilon_{t}^{I}}%
=\beta\mathbb{E}_{t}\left\{  \varepsilon_{t+1}^{C}\left[  \overset{}{\left(
1-\delta\right)  }\frac{\lambda_{t+1}^{c}}{\varepsilon_{t+1}^{I}}%
+\lambda_{t+1}^{c}\textcolor{red}{u_{jt+1}}Z_{t+1}-\lambda_{t+1}%
^{c}\textcolor{red}{\Phi(u_{jt+1})}\right]  \right\} \\
&  :\varepsilon_{t}^{C}\frac{\lambda_{t}^{c}}{\varepsilon_{t}^{I}}%
=\beta\mathbb{E}_{t}\left\{  \varepsilon_{t+1}^{C}\lambda_{t+1}^{c}\left[
\frac{\overset{}{\left(  1-\delta\right)  }}{\varepsilon_{t+1}^{I}%
}+\textcolor{red}{u_{jt+1}}Z_{t+1}-\textcolor{red}{\Phi(u_{jt+1})}\right]
\right\} \\
&  :\frac{1}{\varepsilon_{t}^{I}}=\beta\mathbb{E}_{t}\left\{  \frac
{\varepsilon_{t+1}^{C}}{\varepsilon_{t}^{C}}\frac{\lambda_{t+1}^{c}}%
{\lambda_{t}^{c}}\left[  \frac{\overset{}{\left(  1-\delta\right)  }%
}{\varepsilon_{t+1}^{I}}+\textcolor{red}{u_{jt+1}}Z_{t+1}%
-\textcolor{red}{\Phi(u_{jt+1})}\right]  \right\}
\end{align*}


In a second step, we use the first order condition $\left(  \partial
B_{jt}\right)  $ to remove the Lagrangian multiplier on the budget constraint
$\lambda_{t}^{c}$:%
\begin{align*}
\left(  \partial K_{jt}\right)   &  :\frac{1}{\varepsilon_{t}^{I}}%
=\beta\mathbb{E}_{t}\left\{  \frac{\varepsilon_{t+1}^{C}}{\varepsilon_{t}^{C}%
}\frac{\lambda_{t+1}^{c}}{\lambda_{t}^{c}}\left[  \frac{\overset{}{\left(
1-\delta\right)  }}{\varepsilon_{t+1}^{I}}+\textcolor{red}{u_{jt+1}}Z_{t+1}%
-\textcolor{red}{\Phi(u_{jt+1})}\right]  \right\} \\
&  :\frac{1}{\varepsilon_{t}^{I}}=\frac{1}{R_{t}}\mathbb{E}_{t}\left\{
\left[  \frac{\overset{}{\left(  1-\delta\right)  }}{\varepsilon_{t+1}^{I}%
}+\textcolor{red}{u_{jt+1}}Z_{t+1}-\textcolor{red}{\Phi(u_{jt+1})}\right]
\right\} \\
&  :\frac{R_{t}}{\varepsilon_{t}^{I}}=\mathbb{E}_{t}\left\{  \left[
\frac{\overset{}{\left(  1-\delta\right)  }}{\varepsilon_{t+1}^{I}%
}+\textcolor{red}{u_{jt+1}}Z_{t+1}-\textcolor{red}{\Phi(u_{jt+1})}\right]
\right\}
\end{align*}


Regarding the first order condition on capital utilization, we have:%
\begin{align*}
\left(  \partial u_{jt}\right)   &  :\varepsilon_{t}^{C}\lambda_{t}^{c}\left[
Z_{t}K_{jt-1}-\Phi^{\prime}\left(  u_{jt}\right)  K_{jt-1}\right]  =0\\
\left(  \partial u_{jt}\right)   &  :Z_{t}-\Phi^{\prime}\left(  u_{jt}\right)
=0\\
\left(  \partial u_{jt}\right)   &  :Z_{t}=\Phi^{\prime}\left(  u_{jt}\right)
\\
\left(  \partial u_{jt}\right)   &  :Z_{t}=\bar{Z}+\frac{\psi}{1-\psi}\bar
{Z}\left(  u_{jt}-1\right)
\end{align*}
as $\Phi^{\prime}\left(  u_{jt}\right)  =\bar{Z}+\frac{\psi}{1-\psi}\bar
{Z}\left(  u_{jt}-1\right)  $.

\subsection{Firms}

The modifications are minor here to introduce variable capital
utilization.\ We have to substitute $K_{it-1}$\ by $K_{it}^{u}$ and create a
new variable determining $K_{it}^{u}$:%
\[
K_{it}^{u}=u_{t}K_{it-1}%
\]
Equations where $K_{it-1}$\ was substituted by $K_{it}^{u}$ are given by:
\begin{align}
Y_{it}  &  =\varepsilon_{t}^{A}\left(  \textcolor{red}{K_{it}^{u}}\right)
^{\alpha}H_{it}{}^{1-\alpha},\\
Z_{t}  &  =\alpha\frac{Y_{it}}{\textcolor{red}{K_{it}^{u}}}%
\end{align}


\subsection{Aggregation}

The resource constraint is also affected:%
\[
Y_{t}=C_{t}+I_{t}+\bar{G}\varepsilon_{t}^{G}%
+\textcolor{red}{\Phi(u_{jt})K_{t-1}}
\]


\section{Model summary}

We have to include three new variables
$\{\textcolor{red}{u_{t}},\textcolor{red}{\Phi(u_{t})},\textcolor{red}{K_{t}^{u}}\}$%
, and three corresponding equations to obtain a determinate system of equations:%

\begin{align*}
&  :\beta R_{t}=\mathbb{E}_{t}\left\{  \frac{\varepsilon_{t}^{C}}%
{\varepsilon_{t+1}^{C}}\left(  \frac{C_{jt+1}}{C_{jt}}\right)  ^{\sigma_{C}%
}\right\} \\
&  :W_{t}=\chi H_{t}^{\sigma_{L}}C_{t}^{\sigma_{C}}\\
&  :\frac{R_{t}}{\varepsilon_{t}^{I}}=\mathbb{E}_{t}\left\{  \frac{\left(
1-\delta\right)  }{\varepsilon_{t+1}^{I}}+Z_{t+1}\right\} \\
&  :K_{t}=\varepsilon_{t}^{I}I_{t}+\left(  1-\delta\right)  K_{t-1}\\
&  :\textcolor{red}{K_{t}^{u}=u_{t}K_{t-1}}\\
&  :\textcolor{red}{Z_{t}=\bar{Z}+(\psi/(1-\psi))\bar{Z}(u_{t}-1)}\\
&  :Y_{t}=\varepsilon_{t}^{A}\left(  \textcolor{red}{K_{t}^{u}}\right)
^{\alpha}H_{t}{}^{1-\alpha},\\
&  :Z_{t}=\alpha\frac{Y_{t}}{\textcolor{red}{K_{t}^{u}}}\\
&  :W_{t}=\left(  1-\alpha\right)  \frac{Y_{t}}{H_{t}}\\
&  :Y_{t}=C_{t}+I_{t}+g^{y}\bar{Y}\varepsilon_{t}^{G}%
+\textcolor{red}{\Phi(u_{jt})K_{t-1}}\\
&
:\textcolor{red}{\Phi(u_{t})}=\textcolor{red}{\bar{Z}(u_{t}-1)+(1/2)(\psi/(1-\psi))\bar{Z}(u_{t}-1)^2}\\
&  :\log\left(  \varepsilon_{t}^{A}\right)  =\rho_{A}\log\left(
\varepsilon_{t-1}^{A}\right)  +\eta_{t}^{A}\\
&  :\log\left(  \varepsilon_{t}^{G}\right)  =\rho_{G}\log\left(
\varepsilon_{t-1}^{G}\right)  +\eta_{t}^{G}\\
&  :\log\left(  \varepsilon_{t}^{C}\right)  =\rho_{C}\log\left(
\varepsilon_{t-1}^{C}\right)  +\eta_{t}^{C}\\
&  :\log\left(  \varepsilon_{t}^{I}\right)  =\rho_{I}\log\left(
\varepsilon_{t-1}^{I}\right)  +\eta_{t}^{I}%
\end{align*}


\section{Steady state}

The usual assumption regarding variable capital utilization is to normalize
the steady state rate of utilization $\bar{u}=1$, the steady state for
$\{\textcolor{red}{u_{t}},\textcolor{red}{\Phi(u_{t})},\textcolor{red}{K_{t}^{u}}\}$
boils down to:%
\begin{align*}
\bar{u}  &  =1\\
\Phi(\bar{u})  &  =0\\
\bar{K}^{u}  &  =\bar{K}%
\end{align*}


\section{Dynare writing}

\begin{enumerate}
\item Declare three new endogenous variables
$\{\textcolor{red}{u_{t}},\textcolor{red}{\Phi(u_{t})},\textcolor{red}{K_{t}^{u}}\}$%


\item Declare the elasticity of the utilization rate $\psi$\ as a parameter
and calibrate it to $\psi$\ =0.7

\item Include the three dynamic equations defining
$\{\textcolor{red}{u_{t}},\textcolor{red}{\Phi(u_{t})},\textcolor{red}{K_{t}^{u}}\}$
and modify equations affected by capital utilization.

\item Define the steady state of
$\{\textcolor{red}{u_{t}},\textcolor{red}{\Phi(u_{t})},\textcolor{red}{K_{t}^{u}}\}$%
\ in the init block.
\end{enumerate}


\end{document}